\section{C. Equivariant comparison, drilling, and the topology of special groups}
\label{sec:part-c}

Throughout this part, $U=X\setminus S$ is connected, $X$ is a closed
orientable surface, and $S$ is a finite set of cardinality $\nu$. We put
$A=\operatorname{Homeo}(X,S)$ with its compact-open topology and
$K=\pi_0(A)$. Thus $A$ permits both permutations of $S$ and orientation
reversal; $A^0$ fixes $S$ pointwise and preserves orientation. We call the
type hyperbolic when $2g-2+\nu>0$. A superscript $u$ means the centralizer
of the image of the specified homomorphism $u:\Gamma\to A$. In particular,
$(A^0)^u=A^0\cap A^u$; it is not, by definition, the identity component
of $A^u$. We distinguish these two groups until connectedness is proved.

\subsection{The categories and their exact fidelity conditions}

\begin{claim}{C.1}{Refuted}{The unrestricted full target has too many arrows}
The literal comparison from surfaces and homeomorphisms to the full
$2$-category of peripheral groupoid diagrams and arbitrary functors is
not locally essentially surjective, already for $\Gamma=1$ and type
$(0,3)$. The correct target for a realization biequivalence is the
$2$-groupoid of such diagrams with equivalences as $1$-arrows.
\end{claim}
\begin{proof}
For the three-punctured sphere put $\pi=\langle x,y\rangle$ and take the
peripheral diagram
\[
 \coprod_{i=1}^3 B\Z\longrightarrow B\pi,
 \qquad 1\longmapsto x,\ y,\ (xy)^{-1}.
\]
The identity on its three-object index set, the trivial homomorphism on
each copy of $\Z$, and the trivial homomorphism $\pi\to\pi$ define a
strictly commuting endomorphism of this diagram. Its map on
$\pi_{\rm ab}=\Z^2$ is zero. A homeomorphism induces an automorphism on
this abelianization. Natural isomorphisms between functors induce inner
conjugations, which act trivially on abelianization. Consequently this
endomorphism cannot become isomorphic to the image of a homeomorphism.
This is a counterexample to the literal fullness assertion, not to its
core version.

Here is the precise core version we use. An admissible surface is a
finite disjoint union of finite-type orientable surfaces, each genus-zero
component having at least one puncture, as on scan 264, folio 133. The
source $2$-groupoid has strict $\Gamma$-actions as objects, equivariant
homeomorphisms as $1$-arrows, and endpoint-preserving homotopy classes of
equivariant isotopies as $2$-arrows. An object of the target is a diagram
of groupoids over $B\Gamma$
\[
 D_\Gamma\longrightarrow U_\Gamma\longrightarrow B\Gamma.
\]
Its homotopy fibre over the base object is a peripheral diagram
$\coprod_{i\in I}B\Z\to B\pi$ with the surface presentation
$\prod_{j=1}^{g}[a_j,b_j]\prod_{i=1}^{\nu}c_i=1$; the $i$th map sends
$1$ to $c_i$. The fibre condition includes the common orientation module
and permits the noninjective map in type $(0,1)$ and the repeated image
in type $(0,2)$. This is the condition $M_{\rm lac}$. For $I=\varnothing$
we use the closed surface group, with $g\geq1$.

A target $1$-arrow is a pair of groupoid equivalences over $B\Gamma$
together with the invertible natural transformation expressing
commutativity of the peripheral square. A $2$-arrow is a pair of natural
isomorphisms, with identity projection to $B\Gamma$, satisfying that
square's compatibility. Composition is composition of functors with
pasting of these specified natural transformations. Equivalently these
are equivalences of the corresponding opfibrations over
$[1]\times B\Gamma$ that preserve the displayed fibres. The topological
comparison takes the fundamental groupoids of the homotopy quotients of
a $\Gamma$-invariant punctured collar system and of $U$. Two collar
systems have a common smaller system; the induced equivalences and
natural transformations specify the independence of this choice.
No arbitrary functor is silently replaced by an equivalence.
\end{proof}

\begin{claim}{C.2}{Proved}{The coset category and the quotient correction}
For a hyperbolic peripheral pair $(\pi;(L_i)_{i\in I})$, the category
$\mathscr C$ with objects $u,p_i$, automorphism groups
$\Aut(u)=\pi$, $\Aut(p_i)=1$, and
$\Hom(p_i,u)=\pi/L_i$ determines the pair up to an isomorphism and
simultaneous inner conjugation. Its equivalence groupoid has components
$\Out_{\rm lac}(\pi)$ and trivial automorphism groups. For an action with
extension $1\to\pi\to E\to\Gamma\to1$, the corresponding category over
$B\Gamma$ uses the bimodule $E/L_i$, with right action of
$E_i/L_i=\Gamma_i$, not $E/E_i$.
\end{claim}
\begin{proof}
The stabilizer in $\Aut(u)$ of a selected arrow $p_i\to u$ is $L_i$.
Changing that arrow to $aL_i$ conjugates its stabilizer by $a$. An
equivalence fixes the unique object with nontrivial noncyclic
endomorphism group and permutes the other objects. Its restriction at
$u$ is an automorphism $\alpha$ of $\pi$. Equivariance of its maps on
Hom-sets says precisely that $\alpha(L_i)$ is conjugate to
$L_{\sigma(i)}$. Such a map of transitive sets is determined by a coset
whose stabilizer has this property. The equality
$N_\pi(L_i)=L_i$ makes that coset unique. A natural transformation between
two such equivalences has its component $a\in\pi$ at $u$, its identity
components at the $p_i$, and imposes $\alpha'=\operatorname{Int}(a)\alpha$
and the corresponding multiplication of all cosets by $a$. An
automorphism of an equivalence therefore lies in $Z(\pi)=1$.
This proves the assertion about the equivalence groupoid and gives both
comparison constructions explicitly.

For the geometric exit interpretation, the space of exits from $p_i$ to
$u$ is the homotopy fibre of the link map $BL_i\to B\pi$. Its components
are $\pi/L_i$ and its fundamental groups are the kernels of
$L_i\to\pi$; its higher homotopy groups vanish. The peripheral injection
makes those components contractible. Thus the ordinary category above
has exactly these exit mapping spaces. This also identifies the exact
failure of this ordinary model for the plane: the fibre of
$B\Z\to *$ is $B\Z$, not a discrete space. More explicitly the
Grothendieck category of $BL_i\to B\pi$ has automorphism groups
$L_i,\pi$ and forward Hom-set $\pi$. Its comparison with the exit
category is identity at the generic group, trivial at the special
endomorphism group, and $\pi\to\pi/L_i$ on forward arrows. Conversely
the translation groupoid of the generic action on the exit Hom-set
is $(\pi/L_i)//\pi\simeq BL_i$, with its map to $B\pi$. Thus the
peripheral diagram is recovered by this explicit link construction;
these two displayed categories are not silently identified. On
classifying spaces the forward comparison collapses each link to its
special point, realizing the filling map
$B\pi\to B\pi\cup_{\coprod BL_i}\coprod *$.

In the action category take one special object for each $\Gamma$-orbit
of punctures, choosing $i$ in that orbit. The generic object has
automorphism group $E$, the special object has automorphism group
$\Gamma_i$, and composition
on $E/L_i$ is left multiplication by $E$ and right multiplication by
$E_i/L_i$. These actions commute. The homotopy fibre over the base object of
$B\Gamma$ has special objects indexed by $\Gamma/\Gamma_i$ for each
chosen representative, and hence recovers every upstairs puncture. An
equivalent unskeletonized model uses all punctures and the nonempty
$\Gamma$-transporter sets between those in the same orbit. The
projection of a morphism $eL_i$
to $B\Gamma$ is the image $\bar e\in\Gamma$; it is well-defined because
$L_i\subset\pi$. It is not well-defined on $E/E_i$ when
$\Gamma_i\ne1$: representatives $e$ and $ea$, $a\in E_i$, can have
different images in $\Gamma$.

For a geometric example take $U=\mathbf P^1\setminus\{0,\infty,1,-1\}$,
$\Gamma=C_2$ acting by $z\mapsto-z$, and quotient coordinate $w=z^2$.
Then $E=\pi_1(\mathbf P^1\setminus\{0,1,\infty\})=F_2$. With $a$ a
loop about $w=0$, the upstairs inertia is $L_0=\langle a^2\rangle$
and $E_0=\langle a\rangle$. The right $C_2$ action on
$E/\langle a^2\rangle$ is nontrivial, whereas replacing it by
$E/\langle a\rangle$ erases that action. This is the exact morphism
obstruction to the quotient-category inference in the prompt.
\end{proof}

\begin{claim}{C.3}{Proved}{The four levels and the coefficient ladder}
Let $A$ be a locally path-connected topological group whose components
are $1$-types. For the comparison from strict actions to weak actions in
the categorical group $\Pi_1A$, the first three fidelity conditions are,
on each comparison $M\to N$ of transport spaces,
\[
\begin{array}{c|l}
0&\pi_1M\to\pi_1N\text{ injective at every point},\\
1&\pi_0M\to\pi_0N\text{ injective and }\pi_1M\to\pi_1N\text{ isomorphic},\\
2&\pi_0M\to\pi_0N\text{ bijective and }\pi_1M\to\pi_1N\text{ isomorphic}.
\end{array}
\]
Level $3$ adds essential surjectivity on action objects. When $A^0$ is
simply connected, these give exactly scans 279--280: vanishing of
$\pi_1((A^0)^u)$; connectedness of $(A^0)^u$; conjugacy of two lifts of the
same $\Gamma\to K$ by an element of $A^0$; and existence of a lift.
\end{claim}
\begin{proof}
The source Hom-groupoid is $\Pi_1M$. A functor of groupoids is faithful
exactly when it is injective on the automorphism group at each object.
It is fully faithful exactly when that condition is an isomorphism and
two source objects whose images are connected are already connected.
It is an equivalence exactly when it is fully faithful and every target
component is hit. Applying these three elementary tests proves the
table, including the cumulative meaning of each level.

When $\pi_1(A^0)=0$, the categorical group $\Pi_1A$ is equivalent to the
discrete group $K$. For strict homomorphisms $u,v:\Gamma\to A$ the map
of transport spaces is
\[
 \{a\in A:a u(\gamma)a^{-1}=v(\gamma)\ (\gamma\in\Gamma)\}
 \longrightarrow
 \{k\in K:k\bar u(\gamma)k^{-1}=\bar v(\gamma)\}.
\]
A nonempty source is a torsor under $A^u$. Its fibres over points of the
target are torsors under $(A^0)^u$. The first two lines of the table are
therefore the first two conditions stated. For the third, given a target
transporter $k$, choose $a\in A$ representing $k$. The two actions
$a u a^{-1}$ and $v$ have the same projection to $K$; conjugating them
by an element of $A^0$ produces the prescribed transporter. Conversely,
apply transporter surjectivity to the identity element of $K$. This
also verifies the stronger marginal correction on scan 280.

For comparison with the GALOIS ladder, first take the $1$-truncations
of the transport spaces; fixed spaces are not assumed to inherit the
ambient group's higher truncation. On a map of $1$-types $M\to N$, set-valued constant coefficients test bijectivity on $\pi_0$;
full faithfulness of restriction tests that bijectivity together with
surjectivity on $\pi_1$; equivalence tests bijectivity on $\pi_0$ and an
isomorphism on $\pi_1$. For groupoid-valued coefficients all three tests
coincide with a weak equivalence, because $\pi_iM=\pi_iN=0$ for $i\geq2$;
the same is true for $2$-type coefficients. These assertions follow by
applying the long exact sequence to the homotopy fibres: an $r$-connected
fibre gives isomorphisms through degree $r$ and surjectivity in degree
$r+1$. Thus level $2$ is the $r=0$ equivalence test on the transport
spaces. Levels $0$ and $1$ do not equal a rung in that ladder: their
$\pi_1$-injectivity and $\pi_0$-injectivity tests have been written out
above, whereas each rung of the coefficient ladder demands a component
bijection. Level $3$ concerns objects rather than this map of Hom-spaces.
This is a comparison of the specified maps, not an identification based
only on the words ``faithful'' and ``equivalence''.
\end{proof}

\begin{claim}{C.4}{Proved}{Coherent actions and rectification on spaces}
For a discrete group $P$, the space of self-homotopy equivalences of
$BP$ has components indexed by $\Out(P)$, each with homotopy type
$BZ(P)$. In particular, for the centre-free surface groups the map from
outer actions to coherent actions is an equivalence of action groupoids.
Every coherent action of a discrete group on a space has a strict
realization on a space of the same homotopy type. For $P=\Z$ or $\Z^2$,
the obstruction class in degree $3$ is zero and the coherent actions
with a specified linear action are classified by $H^2(\Gamma,P)$.
\end{claim}
\begin{proof}
The groupoid of functors $BP\to BP$ has homomorphisms $P\to P$ as
objects and conjugating elements of $P$ as arrows. Its equivalence
subgroupoid has automorphisms as objects. The stabilizer of an
automorphism is $Z(P)$, and the orbits are $\Out(P)$. The nerve computes
the mapping space because nerves of groupoids represent $1$-types.
Delooping the resulting categorical group has $\pi_1=\Out(P)$,
$\pi_2=Z(P)$ and no higher homotopy groups. Its Postnikov class is the
usual extension obstruction in $H^3(\Out(P);Z(P))$. For $P$ abelian the
map $\Aut(P)\to\Out(P)$ is the identity and gives a strict section;
the categorical group is $BP\rtimes\Aut(P)$, so this Postnikov class
vanishes. A choice of coherence is then a $2$-cocycle with values in the
specified $\Gamma$-module $P$; changing it by a coherent equivalence
changes the cocycle by a coboundary. This proves the stated
classification without asserting a nonzero obstruction for either
abelian surface group.

For rectification, represent a coherent action by a fibration
$p:Y\to B\Gamma$ with the prescribed homotopy fibre. Pull it back along
$E\Gamma\to B\Gamma$. The pullback $Y'=E\Gamma\times_{B\Gamma}Y$ has a
strict deck action of $\Gamma$. Its projection to the contractible
space $E\Gamma$ has fibre the original homotopy type, so the inclusion
of any chosen fibre is a homotopy equivalence. Its homotopy quotient
recovers $Y\to B\Gamma$. This construction applies simultaneously to
peripheral diagrams by the diagrammatic pullback. It produces a space
model; the additional demand that the model be the original surface is
exactly the realization problem addressed below. For continuous
profinite actions the objects and the homotopy quotient must instead be
specified in a continuous/profinite homotopy theory; the preceding
construction makes no untyped transfer to that setting.
\end{proof}

\subsection{Finite actions on hyperbolic surfaces}

\begin{claim}{C.5}{Proved}{Existence, with punctures and orientation reversal}
For a hyperbolic type and every finite group $\Gamma$, each
homomorphism $\rho:\Gamma\to K$ lifts to a homomorphism
$u:\Gamma\to A$. Its image can be chosen to consist of isometries of a
complete finite-area hyperbolic metric on $U$.
\end{claim}
\begin{proof}
Let $H=\rho(\Gamma)$, a finite subgroup of the extended mapping class
group. Kerckhoff's Theorem 5, together with the explicit extension in
\S IV, pp.~255--256, of the original paper \cite{C_Kerckhoff1983}, gives
a fixed hyperbolic structure for $H$. That extension explicitly permits
cusps; the paper's mapping-class convention explicitly includes
orientation-reversing maps. We use the original paper, rather than
inferring the punctured case from the announcement.

Choose the isometry in each class of $H$. It is unique. Indeed an
isometry isotopic to the identity has a lift to $\mathbb H^2$ commuting
with the surface deck group. That deck group contains two hyperbolic
elements with different endpoint pairs, so an isometry commuting with
both fixes at least three points of the ideal boundary and is the
identity. The same argument covers a finite-area punctured surface;
its non-elementary deck group again contains those elements. Products
of the chosen isometries represent products in $H$, and uniqueness
therefore makes the choices a homomorphism $H\to A$. Composing with
$\rho$ proves the claim. The isometries extend to the end compactification
and permute $S$, so the target really is the specified group $A$.
\end{proof}

\begin{claim}{C.6}{Proved}{Uniqueness up to the identity component}
Let $u,v:\Gamma\to A$ be homomorphisms with $\Gamma$ finite, $U$
hyperbolic, and $\bar u=\bar v:\Gamma\to K$. There is an
$h\in A^0$ such that $h u(\gamma)h^{-1}=v(\gamma)$ for every $\gamma$.
This holds for topological actions without a smoothness assumption.
\end{claim}
\begin{proof}
We first justify the analytic model for either action. For a point
$x\in U$ and its finite stabilizer $H$, lift each element of $H$ to the
universal cover with the chosen lift of $x$ fixed. Uniqueness of a lift
with a specified value makes these lifts a finite group. Compactifying
$\widetilde U\simeq\R^2$ by one point gives a finite action on $S^2$
with two common fixed points. Kolev's compact-sphere theorem
\cite[Theorem 1.2]{C_Kolev2006} conjugates that action into $O(3)$.
A subgroup of $O(3)$ fixing two distinct points is contained in the
orthogonal group of the perpendicular plane; choosing a sufficiently
small round disc at the first point and pulling it back yields a common
invariant disc with an $O(2)$ action. Choosing it small enough makes its
translates by distinct cosets of $H$ disjoint. At a puncture we apply the
same finite-stabilizer disc argument on the compactification, as in the
finite-disc theorem proved from the same source.

The quotient consequently has ordinary charts, cone charts, mirror
charts, and corner charts with their finite cyclic or dihedral local
groups. A finite triangulation respecting these strata supplies a
smooth orbifold structure: smooth each triangle, use collar charts
along its edges, and use the displayed linear quotient charts at its
vertices. Pulling the charts to $X$ makes the original finite action
smooth for a compatible smooth structure. Average a metric and use its
conformal class. Uniformization of the punctured conformal surface gives
a finite-type structure for which each operator is conformal or
anticonformal. Thus we have conformal structures $c_u,c_v$ on the same
marked topological surface.

The extremal quasiconformal map $h:(U,c_u)\to(U,c_v)$ in the identity
isotopy class is unique, by the finite-type Teichm\"uller theorem
\cite[\S3(A)--(D)]{C_Bers1981}. For each $\gamma$, the map
$v(\gamma)^{-1}h u(\gamma)$ is orientation preserving: $u(\gamma)$ and
$v(\gamma)$ have the same orientation sign. Pre- and postcomposition by
these conformal or anticonformal maps leaves the maximal dilatation
unchanged. The equality of the two mapping-class homomorphisms says
that this map is isotopic to the identity, with the punctures retained.
Uniqueness therefore gives $v(\gamma)^{-1}h u(\gamma)=h$. The
quasiconformal map extends across the finite punctures; hence
$h\in A^0$ and has the asserted conjugacy property.
\end{proof}

\begin{claim}{C.7}{Proved}{The finite centralizer component, including mirrors}
For a finite action $u$ on a hyperbolic $U$, the group $(A^0)^u$ is
contractible. In particular it is the identity component of $A^u$.
\end{claim}
\begin{proof}
Replace $\Gamma$ by its image $H$. The preceding local construction
makes $H$ a group of hyperbolic isometries. Let $\pi$ be the surface deck
group and let $E$ be the full group of lifts of $H$ to $\mathbb H^2$.
It is a discrete group of isometries, possibly with reflections; its
quotient $Q$ is a finite-type hyperbolic orbifold. Write $D$ for descent
of a centralizing homeomorphism to the labelled underlying surface of
$Q$, with cone points, corner points, mirror boundary, and punctures
retained.

We give the injectivity argument at the place where an
orientation-preserving-only Birman--Hilden theorem would not suffice.
For $f\in(A^0)^u$, the identity isotopy supplies a lift $\widetilde f$
that commutes with $\pi$. Since $f$ commutes with $H$, conjugation by
$\widetilde f$ induces an automorphism $\alpha$ of $E$ that is the
identity on $\pi$. For $e\in E$ and $a\in\pi$,
\[
 \alpha(e)a\alpha(e)^{-1}=\alpha(eae^{-1})=eae^{-1}.
\]
Thus $e^{-1}\alpha(e)$ centralizes $\pi$, and hence is the identity by
the distinct-endpoint argument in C.5. Consequently $\alpha=\id_E$.
The orbifold Dehn--Nielsen--Baer injectivity theorem, with its mirror
and corner structure, makes $D(f)$ isotopic to the identity through
homeomorphisms preserving those structures; the relevant precise form
is \cite[Theorem 2 and Proposition 3]{C_KolbeEvans2021}. At the algebraic
step in their Lemma 3 we have used the calculation just written:
centralizers of individual reflections need not be cyclic, but the
centralizer of the non-elementary surface subgroup is trivial. Thus
that inapplicable sentence of their proof is not being used. There is an orientation hypothesis in that source when the labelled
underlying quotient is orientable, and we verify it here. The canonical
$\pi$-equivariant lift $\widetilde f$ extends by the identity on
$\partial_\infty\mathbb H^2$. For a compact surface this follows from
bounded displacement on a compact fundamental domain. For a punctured
surface use a compact truncation and disjoint invariant horodiscs:
displacement is bounded on the truncated part, and preservation of each
end sends sufficiently deep cusp horodiscs into the corresponding cusp
horodiscs. A sequence escaping through distinct horodiscs has the same
ideal limit before and after the map, since their visual diameters tend
to zero; a sequence in one horodisc has its cusp endpoint as its only
ideal limit. These observations prove the claimed continuous extension.
Because $\alpha=\id$, $\widetilde f$ preserves each reflection axis.
It fixes that axis's two ideal endpoints individually, so preserves its
linear orientation, and, being orientation preserving on $\mathbb H^2$,
preserves its two sides. It therefore preserves every mirror chamber.
On such a chamber the quotient orientation is the one induced from
$\mathbb H^2$; hence $D(f)$ preserves it. Without mirrors an orientable
quotient is covered orientably and the same assertion is immediate.
Thus the theorem is applied inside its stated homeomorphism group.

Let $\mathcal H(Q)^0$ denote the identity component of that labelled
homeomorphism group. Every path starting at its identity lifts uniquely,
starting at the identity upstairs, to a path centralizing $H$. Lift
first on the regular covering locus. The local cyclic and dihedral
models extend the lift over cone points, corners and mirrors; continuity
of the extended lift follows in the charts from the finite quotient
map. The same construction for a parameter family proves the path and
homotopy lifting properties. The only ambiguity for an endpoint is a
central deck transformation. Such a deck transformation cannot belong
to $A^0$ unless it is the identity, by C.5's isometry argument. It follows
that descent identifies $(A^0)^u$ with $\mathcal H(Q)^0$, as soon as the
latter is simply connected; the construction is continuous in local
isotopy charts, so this is an isomorphism of topological groups.

We check that this group is contractible, including the small underlying
quotients. End-compactify the quotient and retain its finitely many
interior and boundary marked points. Yagasaki's compact-pair Lemma 2.1
\cite{C_Yagasaki2000} gives contractibility outside its explicit
exceptional pairs. A marked point on a boundary component is covered
by Lemma 2.1(ii). With no such point, the exceptional disk has at most
one interior marked point, the annulus and M\"obius band have none, the
sphere has at most two, the projective plane at most one, and the torus
and Klein bottle have none. Each of these has nonnegative orbifold
Euler characteristic: an interior cone subtracts $1-1/m$, an interior
puncture subtracts $1$, and a corner or boundary puncture would already
be a boundary marked point. But
$\chi_{\rm orb}(Q)=\chi(U)/|H|<0$. None occurs. The homeomorphism group
with labels has the same identity component as the group fixing all
these finite marks pointwise. Yagasaki's theorem therefore applies.
This proves contractibility, and also supplies the simple connectedness
used in the lifting construction.
\end{proof}

\begin{claim}{C.8}{Proved}{The finite core biequivalence and homotopy fixed points}
On hyperbolic surfaces, for every finite $\Gamma$ the corrected
comparison of C.1 has all four fidelity levels and is a biequivalence.
For a specified action $u$, the comparison
\[
 A^u\longrightarrow A^{h\Gamma}
      =\operatorname{Map}_\Gamma(E\Gamma,A),
 \qquad a\longmapsto(\text{constant map with value }a),
\]
where $\Gamma$ acts on $A$ by conjugation through $u$, is a weak homotopy
equivalence. In standard CW models it is a homotopy equivalence.
\end{claim}
\begin{proof}
The nonequivariant peripheral comparison identifies $K$ with
$\Out_{\rm lac}(\pi)$ and has trivial $2$-automorphism groups, by the
peripheral Dehn--Nielsen--Baer theorem
\cite[Theorems 8.1 and 8.8]{C_FarbMargalit2012} and $Z(\pi)=1$.
C.7 gives levels $0$ and $1$, C.6 gives level $2$ by the transporter
argument of C.3, and C.5 gives level $3$. These arguments also treat a
nonfaithful homomorphism: its finite image is the action group, and a
finite-order homeomorphism in $A^0$ is the identity. Disconnected
admissible surfaces reduce to representatives of the finitely many
orbits of connected components and their stabilizers, with the stated
transporters providing the arrows between choices. This is the desired
biequivalence on its hyperbolic part.

The map $A\to K$ is $\Gamma$-equivariant. Its fibres are contractible by
the same homeomorphism-group theorem with $H=1$. Taking derived
homotopy fixed points gives a weak equivalence
$A^{h\Gamma}\to K^\Gamma=C_K(\bar u(\Gamma))$, the latter a discrete
space. The constant-map comparison followed by this map is exactly
$\pi_0(A^u)\to C_K(\bar u(\Gamma))$. Its surjectivity follows by applying
C.6 to $a u a^{-1}$ and $u$ for any representative $a$ of a centralizing
mapping class; its injectivity and its contractible fibres are C.7.
Thus the displayed map is a weak equivalence. The source components
have CW homotopy type by the ANR statement for quotient homeomorphism
groups in \cite{C_Yagasaki2000}; taking the usual CW model of the
derived target turns it into a homotopy equivalence. We state the model
convention because scan 281 writes ``homotopisme'' without specifying
an infinite-dimensional mapping-space model.
\end{proof}

\subsection{The exceptional types and infinite operator groups}

\begin{claim}{C.9}{Proved}{Corrected abelian tests and a torus calculation}
For $\pi_1(A^0)\ne0$, replacement of the four tests by connectedness and
uniqueness of lifts of $\Gamma\to K$ loses the coherent action data.
The valid tests are C.3 with weak transport spaces. At a specified
strict action the target automorphism space is the $1$-type
$A^{h\Gamma}$ and
\[
 \pi_1(A^{h\Gamma},1)=\pi_1(A,1)^\Gamma.
\]
For $A=\operatorname{Homeo}(T^2)$ and $u(C_2)=\{1,-1\}$, the group
$(A^0)^u$ has four components, all contractible, whereas the right-hand
invariant fundamental group is zero. Thus the naive connectedness test
fails; the corrected comparison does not fail in this example.
\end{claim}
\begin{proof}
The homotopy-fixed-point groupoid of a $1$-type has loops equal to the
fixed automorphisms of its base object, proving the displayed equality.
Its components include the $H^1(\Gamma;\pi_1A)$ coherence data, with
possible obstruction for an individual component of $A$. Hence it is
not legitimate to test only $\pi_0A$ when $\pi_1A\ne0$.

For the torus, evaluation at $0$ gives
$A^0\cong T^2\times A^0_{\{0\}}$ by
$f\mapsto(f(0),t_{-f(0)}f)$. The second factor is contractible by the
homeomorphism-group theorem for a torus with a marked point. In the
centralizer of $-1$, $f(0)$ is one of the four points of $T^2[2]$.
Each is realized by translation. In the fibre $f(0)=0$, regarding the
torus as once punctured makes it hyperbolic; the fibre is contractible
by C.7. This proves the four-component assertion. Conjugation by $-1$
acts as $-I$ on $\pi_1(A^0)=\Z^2$, so its invariants vanish.
The inclusion of translations $T^2\hookrightarrow A^0$ is equivariant
and a weak equivalence; its homotopy fixed points have components
$H^1(C_2;\Z^2_{-})=(\Z/2)^2$ and zero higher homotopy groups. The four
translations just identified realize exactly these four components.

For the sphere, the peripheral groupoid diagram is empty over $*$ and
cannot encode its homotopy type; $\pi_2(S^2)=\Z$ already records the
failure of asphericity. The plane, annulus and torus have
$A^0\simeq S^1,S^1,T^2$, respectively; in the plane the peripheral
map $B\Z\to *$ retains the circle of automorphism $2$-arrows, and in the
annulus the two ends remain indexed even though their subgroup images
coincide. These exact maps explain why collapsing to a bare outer group
would erase data. The complete finite-action comparison for these models, including
its parametric centralizer proof, is C.10.
\end{proof}

\begin{claim}{C.10}{Proved}{Every finite action in the exceptional aspherical types}
For $U=\R^2$, $S^1\times\R$, or $T^2$, and every finite strict action,
the fixed-to-homotopy-fixed comparison in C.8 is a weak equivalence
and is a homotopy equivalence on CW models. The corrected peripheral
$2$-groupoid comparison has all four fidelity levels for these types.
The tests are those of C.3; ordinary outer actions alone do not retain
the required coherence classes.
\end{claim}
\begin{proof}
We prove the parametric assertion, not merely classification of individual
actions. Write respectively
\[
 G=O(2),\qquad
 G=(\R/\Z)\rtimes(C_2\times C_2),\qquad
 G=(\R^2/\Z^2)\rtimes\operatorname{GL}_2(\Z).
\]
The annulus element $(a,\epsilon,\eta)$ sends
$(\theta,t)$ to $(\epsilon\theta+a,\eta t)$; the first sign acts on the
circle factor and the second on the two ends. In every case $G\to A$
is a weak equivalence. Here is a direct check of the identity components.
On the torus, evaluation at $0$ splits $A^0$ as $T^2$ times its
point-fixing subgroup. An identity isotopy can be made point fixing by
postcomposing at time $s$ with translation by the negative of the trace
of $0$. Thus that subgroup is its own identity component, and is
contractible by the punctured-torus case of the homeomorphism-group
theorem \cite[Theorem 1.1 and Lemma 2.1]{C_Yagasaki2000}. The same adjustment on the annulus uses angular rotation and
translation of its real coordinate. Its point-fixing subgroup is the
identity group of a hyperbolic three-ended sphere and is contractible,
so $A^0\simeq S^1$. On the plane first fix $0$ by translations and
then fix $1$ using multiplication by a nonzero complex number. The
remaining identity group fixes $0,1,\infty$ and is contractible; hence
$A^0\simeq S^1$ with generator the rotation circle. The component
groups are respectively $C_2$, $C_2\times C_2$, and
$\operatorname{GL}_2(\Z)$: orientation and end signs give the first
two, and the torus action on its lattice gives the third. Cutting the
plane or annulus along a ray or spanning arc verifies the stated
kernels; the torus case is the genus-one Dehn--Nielsen--Baer theorem.

By the simultaneous finite normal forms proved in B.1 and B.14,
conjugate the given finite image $H$ into $G$. We will prove
\[
                    C_G(H)\longrightarrow C_A(H)                 \tag{C.10.1}
\]
is a weak equivalence by examining its fibres. Conjugation by the chosen
homeomorphism transports both the fixed and derived-fixed comparisons,
so this reduction is an isomorphism of the actual diagrams.

\emph{Torus displacement.}
For $f\in C_A(H)$ let $M\in\operatorname{GL}_2(\Z)$ be its lattice
action and lift it as $F(x)=Mx+d(x)$, with $d$ periodic. Put
$b=\int_{T^2}d(x)\,dx$ and $g(x)=Mx+b$ modulo $\Z^2$.
Changing the lift changes $b$ by an integral vector and leaves $g$
unchanged. For $\gamma(x)=A_\gamma x+t_\gamma$ the commutation equation
is
\[
 d(A_\gamma x+t_\gamma)-A_\gamma d(x)
       =(I-M)t_\gamma+k_\gamma,\qquad k_\gamma\in\Z^2.
\]
Integration preserves its left-hand change of variables, because the
affine operator preserves Haar measure. It gives the same equation
with $d$ replaced by $b$. Thus $g$ commutes with $H$ and
$g^{-1}f\in A^0$ has a unique lift of zero mean displacement.
Consequently multiplication is a homeomorphism of spaces
\[
 C_G(H)\times\mathcal K\ \xrightarrow{\sim}\ C_A(H),
\]
where $\mathcal K$ consists of these zero-mean maps. The averaging
map need not be a group homomorphism; only this explicit space
factorization is used. A zero-mean lift commutes exactly with every
lift of $H$: the possible integral commutator above has zero average
and is therefore zero.

First divide by the finite translation subgroup $N=H\cap T^2$.
If $a$ represents a member of $N$ of order $m$, commutation implies
$F(x+a)=F(x)+a+k$. Iterating $m$ times and using $ma\in\Z^2$
forces $k=0$ for a zero-mean identity lift. Hence its displacement is
periodic for the enlarged covering lattice. Descent and lift identify
$\mathcal K$ with the same zero-mean space on $T'=T^2/N$ for $H/N$.
This quotient group has faithful linear part, a finite subgroup $P$
of $\operatorname{GL}_2(\Z)$.

When $P$ contains a nonidentity positive operator, its positive subgroup
$K$ is cyclic, generated by an affine rotation with $I-A$ invertible.
It has a nonempty finite fixed set $R$ on $T'$, stable under $H/N$.
A zero-mean lift commutes with each lift of that rotation and fixes its
unique Euclidean fixed point; thus it fixes every point of $R$.
Choose $p\in R$. The point-fixing argument above places our map in
$\operatorname{Homeo}(T',p)^0$. C.7, applied to the cyclic action on
$T'\setminus\{p\}$, supplies an equivariant identity isotopy. This
isotopy fixes all of $R$, since $R$ is finite. Our map therefore belongs
to $\operatorname{Homeo}(T',R)^0$. Apply C.7 again, now to the full
$H/N$ action on the hyperbolic $T'\setminus R$. Its identity centralizer
is exactly $\mathcal K$: conversely an equivariant identity isotopy
relative to $R$ has commuting identity lifts, whose mean is fixed by
$A$ and hence is zero. This proves contractibility of $\mathcal K$.

The remaining linear groups are $P=1$ and a group of order two generated
by a reflection. For $P=1$, the space factorization on the identity
component is $\operatorname{Homeo}(T')^0\cong T'\times\mathcal K$.
The first factor is a section of the mean map. Since the left side has
the homotopy type $T'$, the mean map is an isomorphism on $\pi_1$ and
$\mathcal K$ is weakly contractible. It is a retract of the ANR on the
left, so it has CW homotopy type and is contractible.

For a reflected or glide-reflected torus the underlying quotient is an
annulus, a M\"obius band, or a Klein bottle. This exhausts the cases:
the action reverses orientation, has no isolated rotation point, and
the quotient has Euler characteristic zero; its boundary consists of
mirrors. A zero-mean lift $F$ commutes exactly with the Euclidean orbifold
group $E$. Each lifted reflection line is preserved, and bounded
displacement prevents reversal of its infinite directions. In
particular the compact annulus quotient's two boundary circles cannot
be interchanged. The explicit straight homotopy
\[
                  F_s(x)=(1-s)F(x)+sx
\]
is $E$-equivariant, because every element of $E$ is affine. On each
reflection line it stays on that same line. It therefore descends to
a homotopy of the quotient as a surface with boundary, preserving its
boundary throughout. Its quotient domain is compact, so it is proper.
Epstein's Theorem 6.3, with the topological Appendix
\cite{C_Epstein1966}, applies with exactly this proper,
boundary-preserving homotopy and gives an identity isotopy. This also
handles the boundary-free Klein bottle; Theorem 5.8 of that source
separately verifies the M\"obius-band case. Thus no unverified inference
from an abstract outer automorphism to an exceptional-surface isotopy
is needed here.

Descent consequently puts $\mathcal K$ in the identity group of the
labelled quotient. Every identity isotopy there lifts; among the two
lifts exactly one preserves the orientation upstairs. This identifies
the identity component of $C_{A^0}(H/N)$ with the quotient identity
homeomorphism group. By Yagasaki's actual exception list its homotopy
type is $S^1$ in all three cases. The mean factorization on that component
is $T'^P_0\times\mathcal K$, with $T'^P_0\cong S^1$. Here
$P$ acts linearly on the translation torus; this fixed translation
circle is present even for an affine glide with no fixed point.
The section by translations makes the mean map surjective on $\pi_1$;
a surjection $\Z\to\Z$ is an isomorphism. Hence $\mathcal K$ is
weakly contractible, and is contractible by the same ANR-retract argument.
This completes (C.10.1) for the torus, including free glide reflections.

\emph{Annulus displacement.}
Lift $f\in C_A(H)$ to
$F(\theta,t)=(\epsilon_f\theta+d(\theta,t),F_2(\theta,t))$ and put
$b=\int_0^1d(\theta,0)\,d\theta$. Define
$g(\theta,t)=(\epsilon_f\theta+b,\eta_f t)$, where $\eta_f$ is its
end sign. The commutation equation, integrated on the invariant core
$t=0$, proves that $g$ commutes with $H$. It gives the same continuous
factorization $C_A(H)\cong C_G(H)\times\mathcal K$, with zero-mean
identity lifts in $\mathcal K$. Divide first by the finite rotation
subgroup. The resulting faithful sign group is a subgroup of
$C_2\times C_2$.

When it contains $(-1,-1)$, that positive involution has two fixed
points. Its subgroup is normal. A zero-mean lift fixes each of the
corresponding fixed points in $\R\times\R$. The same two applications
of C.7 as in the torus argument, first marking one fixed point and then
the whole invariant fixed set, identify $\mathcal K$ with a contractible
centralizer on a hyperbolic punctured annulus. An equivariant identity
lift in that punctured model has mean zero because its angular linear
sign is $-1$, proving equality in both directions.

A subgroup without $(-1,-1)$ is trivial, generated by $(-1,1)$, or
generated by $(1,-1)$. In the first case the factorization and
$A^0\simeq S^1$ give contractibility as above. An angular reflection
has a strip as quotient, with its two mirror boundary lines and two
ends. Exact commutation of a zero-mean lift preserves both lines and
both ends. The same straight homotopy in $\R\times\R$ is equivariant
for all the affine covering operators and preserves the boundary
lines. It is proper on the strip quotient: for $t\to+\infty$ both
$t$ and $F_2(\theta,t)$ tend uniformly to $+\infty$ on a compact
angular fundamental interval, and for $t\to-\infty$ both tend to
$-\infty$. Every convex combination has the same uniform property.
Epstein's proper boundary-homotopy theorem therefore places the
descended map in the identity group. That group is contractible by
Yagasaki, so $\mathcal K$ is contractible.
Finally a radial reflection or radial glide reflection has quotient a
half-open annulus or a punctured projective plane, respectively. The
straight homotopy is again equivariant; it preserves the fixed mirror
line in the reflection case and has the just-verified properness at
the quotient's end. Epstein's Theorem 6.3, or Theorem 5.8 for the open
M\"obius band, gives the required identity isotopy. Descent and isotopy
lifting identify the upstairs centralizer identity component with that
group. Yagasaki gives $S^1$ in both cases. The mean factorization has a
section by the fixed angular circle, hence again forces
$\mathcal K$ to be contractible. These cases include every affine
finite action and retain the end permutation throughout.

\emph{Plane centralizers.}
For a cyclic rotation group of order $n>1$, a commuting map fixes $0$.
Descent through $z\mapsto z^n$ is a covering of the orientation-preserving
centralizer onto $\operatorname{Homeo}^+(\R^2,0)$, with deck group
$C_n$. The base is connected with homotopy type $S^1$. A loop of one
rotation downstairs lifts from the identity to the deck rotation by
$2\pi/n$. Its monodromy is therefore onto $C_n$, so the covering
centralizer is connected and the upstairs rotation circle induces the
isomorphism on $\pi_1$; all higher homotopy groups vanish. A negative
map centralizes the rotation only for $n=2$, as its action on the
punctured-plane fundamental group inverts the deck generator. In that
case conjugation supplies the second component. These are precisely
the components and circle types of $C_{O(2)}(C_n)$.

For a reflection, descent has half-plane quotient. A commuting map has
two independent signs: its orientation on the fixed line and its
orientation on the plane. The four diagonal orthogonal maps realize
them. Each fibre is the identity homeomorphism group of the half-plane,
which is contractible by Yagasaki. For a dihedral group with rotation
order $n>2$, a centralizing map preserves plane orientation and each
reflection line. Preservation of cyclic order of their rays permits
only the identity permutation or simultaneous reversal of every ray;
these are realized by $I$ and $-I$. For $n=2$ the two perpendicular
reflection lines permit four independent ray-sign choices, realized by
the four diagonal orthogonal maps. After dividing by the matching map,
descent is to a sector homeomorphism fixing its vertex and both boundary
rays. This group is contractible: the sector is a half-plane with a
marked boundary point, and its identity group is covered by Yagasaki's
boundary-point case. The boundary rays also make the lift unique.
The trivial group case was verified at the start. This proves (C.10.1)
for every finite subgroup of $O(2)$.

It remains to identify the derived comparison, including all components.
The following elementary torus calculation does so for each of our
three $G$. Let $T=V/\Lambda$, with $V$ real and $\Lambda$ a lattice,
and let a finite group act on a $T$-torsor by affine maps. Its translation
part is a cocycle $b\in Z^1(\Gamma,T)$, where $T$ and $\Lambda$
have the specified linear $\Gamma$-module structures. The homotopy-fixed space is
the section space of its Borel fibration. Its obstruction is the class
$\delta[b]\in H^2(\Gamma,\Lambda)$ of the extension obtained by lifting
the affine maps to $V$. Averaging cochains gives
$H^j(\Gamma,V)=0$ for $j>0$, so
$\delta:H^1(\Gamma,T)\to H^2(\Gamma,\Lambda)$ is an isomorphism.
Thus a homotopy fixed point exists exactly when $b$ is a coboundary,
which is exactly existence of an actual fixed point. After choosing one,
\[
 \pi_0(T^{h\Gamma})=H^1(\Gamma,\Lambda),\qquad
 \pi_1(T^{h\Gamma})=\Lambda^\Gamma,
\]
and there are no higher groups. The exact sequence
\[
 0\to\Lambda^\Gamma\to V^\Gamma\to T^\Gamma
       \to H^1(\Gamma,\Lambda)\to0
\]
identifies the actual fixed components and their fundamental groups
with exactly these groups. The constant-map comparison induces these
identifications. Apply this to every invariant component of $G$, each
of which is just such an affine torus torsor. Noninvariant components
contain neither fixed nor homotopy-fixed points. It follows that
$G^\Gamma\to G^{h\Gamma}$ is a weak equivalence. The equivariant weak
equivalence $G\to A$, together with (C.10.1), now proves the claimed
comparison for $A$ by the commutative square of fixed and derived-fixed
spaces.

Finally $BG$ is the split $2$-type with fundamental group
$\pi_0G$, second group $\Lambda=\pi_1G$, and zero Postnikov invariant.
Coherent finite actions with a specified homomorphism to $\pi_0G$ are
therefore classified by $H^2(\Gamma,\Lambda)$. The same connecting
isomorphism realizes each by a $T$-valued cocycle, hence by a strict
homomorphism $\Gamma\to G$. It also identifies coherent equivalences
in the identity component with coboundaries, namely conjugacies by
translations. The fixed-space calculation supplies their full
transport-groupoid comparison; a prescribed nonidentity transporter
component is treated by composing with its representative in $G$.
Thus the corrected four fidelity levels all hold. We finish by checking
that this is the actual peripheral comparison, not just an abstract
replacement of $A$. For the plane the diagram is $B\Z\to *$.
Its equivalence $2$-group has component group $\{\pm1\}$,
$2$-automorphism group $\Z$, and sign action. For the annulus use
$B\Z\amalg B\Z\to B\Z$ with the two maps $+1,-1$.
An equivalence is specified on components by a generic sign $\epsilon$
and an end permutation of sign $\eta$. Its peripheral sign is forced
to be $\epsilon\eta$. The two square conjugators can be removed by
peripheral natural transformations; the surviving automorphism of
such a diagram equivalence is one diagonal $\Z$, on which the generic
sign $\epsilon$ acts. Thus its $2$-group has exactly the two-sign
component group and lattice action of the annulus $G$ displayed above.
For the torus, the empty peripheral diagram over $B\Z^2$ has
component group $\operatorname{GL}_2(\Z)$ and $2$-automorphism group
$\Z^2$ with its standard action. In all three cases strict diagram
functors with zero square-conjugators give a section, so the
Postnikov class is zero. The actual comparison maps a full plane or
annulus rotation to the generator traced by a peripheral or generic
basepoint, and torus translation loops to the two lattice generators.
It is therefore an isomorphism on these specified $\pi_1$ groups,
as well as on components. This proves that the models and comparisons
used throughout the proof are the ones induced by C.1's peripheral
functor. In particular the annulus end permutation has never been
forgotten.
\end{proof}

\begin{claim}{C.11}{Proved}{Infinite cyclic failures at levels one and two}
For every hyperbolic surface containing a disc, there is a strict
$\Z$-action for which $(A^0)^u$ is disconnected. There also are two
strict $\Z$-actions with the same mapping-class homomorphism that are
not conjugate. Thus finite-group hypotheses cannot be removed from
levels $1$ and $2$.
\end{claim}
\begin{proof}
In a coordinate disc choose two disjoint small round discs centered at
$q_\pm=\pm1$. Define $f$ to rotate each small disc by angle
$\sqrt2(1-r/r_0)$ at radius $r<r_0$, and to be the identity elsewhere.
This is a homeomorphism isotopic to the identity through the same
formula with the angle multiplied by a parameter. The two centres are
isolated fixed points; all other fixed points lie outside the interiors
of these discs. At radius $r_0/2$ the rotation angle is
$\sqrt2/2$, whose ratio to $2\pi$ is irrational, so $f$ has infinite
order. Let $h$ rotate the larger disc by angle $\pi$ on a concentric disc
containing the two small discs, tapering continuously to angle zero
before the outer boundary. It is also isotopic to the identity.
It exchanges the two small discs by rigid rotation and therefore
commutes with $f$.

Set $u(1)=f$. Every homeomorphism commuting with $f$ preserves its fixed
set and sends isolated fixed points to isolated fixed points. Along a
path in the centralizer the induced permutation of $\{q_-,q_+\}$ is
constant. The identity fixes both and $h$ exchanges them. They therefore
belong to different components of $(A^0)^u$. This proves the level-$1$
failure by the exact test C.3. The action $u$ and the trivial $\Z$-action
both project to the trivial homomorphism into $K$, but a conjugate of
$f\ne\id$ cannot be the identity. This proves the level-$2$ failure.
The formulas are checked on an explicit grid in the supplied program;
the isolation and irrationality arguments prove the assertions beyond
that grid.
\end{proof}

\begin{claim}{C.12}{Open}{The requested infinite-group zero-level example}
For the admissible surfaces and the exact comparison of C.1--C.3,
we have not produced an infinite operator group whose level-$0$
comparison has a nontrivial kernel, nor proved universal injectivity.
For a hyperbolic surface the required kernel is
$\pi_1(C_{A^0}(u(\Gamma)),1)$, since the target is zero.
The disconnected centralizer in C.11 proves level-$1$ failure.
The enlarged sphere comparison in C.25 has its stated nontrivial
kernel, but the unpunctured sphere is excluded from C.1's object class.
The requested admissible infinite-group zero-level decision remains
unsettled here.
\end{claim}
The attempted construction using a supported radial rotation produces
isolated fixed points and a component obstruction, not a nonzero loop.
Adding finitely many marked fixed points still gives aspherical
configuration spaces and cannot by that argument produce the required
$\pi_2$ of an evaluation base. No implication from the examples for
levels $1$ and $2$ to level $0$ is asserted. The independent infinite-group
non-realization theorem C.23 gives level-$3$ failure and does not close
this missing zero-level test. In particular, the possibility that level
$0$ holds for all operator groups has not been ruled out by this return.

\paragraph{Scope of the remaining infinite-group obligation.}
The source fixed in C.1 is the admissible-surface category of scan~264,
folio~133: every genus-zero compact component meets the finite puncture
set. The unpunctured sphere is therefore absent. The last sentence of
Part~C(b) asks for an infinite-group example at each fidelity level
within this specified comparison. An example on one of its admitted
abelian types would meet that existential request; a hyperbolic example
would settle the more specific question in C.12. The examples C.11 and
C.23 supply the independent obstructions at levels~$1$, $2$, and~$3$.
The level-$0$ example inside the admitted source remains unprovided.

Here is the exact kernel that is required. For an arbitrary discrete
operator group and a strict action $u:\Gamma\to A$, put
$H=u(\Gamma)$ and let $N_u$ be its weak automorphism groupoid in
$\Pi_1A$. At its identity object, a $2$-automorphism is precisely a
loop class in $A$ invariant under conjugation by every $u(\gamma)$.
The comparison is thus the inclusion-induced map
\[
 \pi_1(C_A(H),1)\longrightarrow
 \pi_1(N_u,1)=\pi_1(A,1)^\Gamma.                 \tag{C.12.1}
\]
Pointwise commutation proves that the image lies in these invariants.
Every based loop and every based homotopy in the source lies in $A^0$,
so its source can equally be written $\pi_1(C_{A^0}(H),1)$.
For hyperbolic types the target vanishes, giving exactly C.12's test.
For the plane, annulus, and torus the target is the invariant subgroup
of $\Z$, $\Z$, or $\Z^2$, respectively; the required obstruction is a
nonzero kernel of (C.12.1).

C.39 proves that every class in this kernel represented by a loop
whose values commute pairwise is zero. Its contraction is global on
the plane, including arbitrary changes and accumulation of orbit types,
and its explicitly lifted filling disk descends to every connected
admissible surface. Thus a remaining example must involve a class
without a commuting-valued representative. C.40 constructs such a representative for every loop whose generated
closure, in the actual compact-open topology, is locally compact.
No representative reduction for arbitrary loops is proved here. C.41 proves the exact
inclusion injective for every operator group with a dense orbit,
without a local compactness assumption on either group. C.36--C.37 separately
provide the stated contraction on the closed union of dimension-two
proper-image open orbits for general loops. C.42 identifies the actual
boundary quotient as an open map and proves that every component of
every real lifted fibre is a properly embedded line; it also supplies
an exact example with disconnected whole circle fibres. C.43 proves
that the real lifted fibre is connected and supplies actual product
coordinates. C.44 gives the five downstairs covering forms, with
source a plane or annulus. C.45 gives the exact relation between two
quotients and its three planar coordinate domains. C.46 constructs a
simultaneous intrinsic contraction when the full set of lifted
boundary fibre partitions has exactly two members, retaining every
operator and quotient relation. C.47 proves that three independent
actual quotients force the full compatible group on the source plane
to be locally compact and Lie, with its actual topology. C.48 gives
the simultaneous intrinsic contraction for the full one-foliation
family, including every component of the proper operator group and
the dense lattice case. C.49 identifies the actual closed lifted
generated group in the three-quotient case and contracts each whole
orbit closure, with exact overlap agreement and finite pasting.
C.50 supplies an exact uniform word bound and extends these
contractions through arbitrary accumulation, giving a joint
contraction on $X_*=\overline{U_2\cup U_3}$. C.51 extends
the intrinsic contractions through the entire frontier when
$O/Q_O$ is compact and the full boundary-quotient family has
at most two members, with the exact median boundary value
and finite pasting. For an empty family every frontier point
is then fixed. C.52 supplies exact agreement with the
prescribed exponential values on $X_*$ by a natural hybrid
target rule and diagonal frontier continuity. C.53 enlarges
that region to $X^\dagger$ using all global nonopen quotient
maps, and controls every accumulation of all $H$-translates
of finitely many remaining cocompact representatives.
C.54 proves actual exponential closure contractions whenever
there are two metabelian global quotient images, including
all diagonal affine weights and the product affine group.
C.55 replaces the proper operator image by the full local
centralizer, with direct equivariant-map limits and the
larger prescribed region $X^\ddagger$. C.56 gives a uniform
word bound for all diagonal one-parameter affine groups and
extends their fractions through arbitrary accumulation on
$X^\sharp$, retaining the full-local relative extension.
C.57 determines the exact full local centralizers and
quotient spaces for every full two-partition family.
The remaining requirements concern untreated noncocompact quotients,
arbitrarily many $H$-inequivalent remaining components, and
extension through nonopen regions outside these closed sets.

C.25 proves such a kernel for the separately enlarged
$\Pi_1\operatorname{Homeo}(S^2)$ comparison. Its relationship with the
admitted punctured comparisons can be computed explicitly. Write
$p_\pm$ for the poles of its irrational rotation. In the identity
component its centralizer $C$ fixes both poles, and it is the same
subgroup in the unmarked, once-marked, and twice-marked sphere
homeomorphism groups. Forgetting marks induces
\[
\begin{gathered}
 \pi_1C=\Z
 \xrightarrow{\;1\;}\pi_1\operatorname{Homeo}(S^2,\{p_+,p_-\})=\Z
 \xrightarrow{\;1\;}\pi_1\operatorname{Homeo}(S^2,\{p_+\})=\Z,\\
 \pi_1\operatorname{Homeo}(S^2,\{p_+\})=\Z
 \xrightarrow{\bmod 2}\pi_1\operatorname{Homeo}(S^2)=\Z/2.
\end{gathered}                                           \tag{C.12.2}
\]
The same full axis-rotation loop generates the first three groups by
the plane and annulus comparisons of C.10; its inclusion into $SO(3)$
gives the last map. Conjugation by the irrational rotation is based
homotopic to the identity in each ambient group, so these are also the
weak-transporter target groups in (C.12.1). Thus puncturing one or both
poles removes C.25's kernel. Every nonpolar point has an infinite orbit,
so this action admits no other finite invariant puncture set.

Keeping the peripheral-diagram recipe while adding the sphere gives
yet another explicitly related target: its diagram is
$\varnothing\to *$, whose equivalence $2$-group is terminal. The
comparison factors through the functor
$\Pi_1\operatorname{Homeo}(S^2)\to\operatorname{Eq}(\varnothing\to*)$,
which sends the identity $2$-automorphism group $\Z/2$ to zero. For the
irrational action its fundamental-group map is the composite
$\Z\to\Z/2\to0$; C.25 computes the first arrow. Neither enlargement
changes the admitted object class fixed in C.1. Consequently C.25 is
retained with its exact stated comparison, while the admissible
infinite-group zero-level obligation remains recorded here.

\paragraph{Exact lift comparison for the remaining zero-level question.}
The compact-open topology agrees with the marked compactification
topology used in C.1. More precisely, restriction
$\operatorname{Homeo}(X,S)\to\operatorname{Homeo}(X\setminus S)$
is a topological-group isomorphism. A homeomorphism extends uniquely
to the finite end compactification, giving the algebraic inverse.
For its continuity at $f_0$, choose disjoint connected target end
neighborhoods $W_j$ of small diameter and complementary compact
core $C$. Choose a source core $K$ with
$f_0^{-1}(C)\subset\operatorname{int}K$, and connected source tails
$V_i$ whose $f_0$-images lie in the prescribed $W_j$.
The conditions $f^{-1}(C)\subset\operatorname{int}K$ and
$f(z_i)\in W_j$ for one point $z_i\in V_i$ are open:
inversion is continuous in the compact-open homeomorphism group
of a locally compact, locally connected surface.
Connectedness forces $f(V_i)\subset W_j$, since this image avoids
$C$ and meets $W_j$. This controls the extended maps uniformly on
all tails. Compact-open control on the remaining compact core
controls them on the rest of $X$, proving continuity of the inverse.

Let $U$ be the hyperbolic surface of C.12, let $A^0=\operatorname{Homeo}(U)^0$, let $H$ be the image of the specified operator group, and let $p:\widetilde U\to U$ be the universal covering with deck group $\pi$. Write $\widetilde H$ for the group of all lifts of elements of $H$, regarded as actual homeomorphisms of $\widetilde U$. There is a continuous injective homomorphism
\[
 j:A^0\longrightarrow\operatorname{Homeo}(\widetilde U),
 \qquad p\,j(f)=f\,p,
\]
characterized by $j(f)\delta=\delta j(f)$ for every $\delta\in\pi$. It identifies the precise centralizers:
\[
 j\bigl(C_{A^0}(H)\bigr)
   =j(A^0)\cap C_{\operatorname{Homeo}(\widetilde U)}(\widetilde H).
\]
In particular every based loop in the left-hand centralizer lifts to a closed loop of homeomorphisms commuting with the entire lift group.
\begin{proof}
Lift an isotopy from $\id$ to $f$ starting at $\id_{\widetilde U}$. Conjugation on the discrete deck group is constant during the lift, so its endpoint commutes with $\pi$. Two such endpoints differ by a deck transformation centralizing $\pi$; they are equal because $Z(\pi)=1$ by Part A. Uniqueness also makes $j$ a homomorphism and an injection. For continuity, fix $\widetilde u\in\widetilde U$. The evaluation map $A^0\to U$ lifts continuously to $\widetilde U$, since $A^0$ is contractible and locally path connected by the surface homeomorphism-group theorem. This lift gives $j(f)(\widetilde u)$, and the local lifting property of a covering gives continuous dependence of the lifted homeomorphism in the compact-open topology.

Suppose $f$ centralizes $H$ and $\widetilde h$ is any lift of $h\in H$. Then $[j(f),\widetilde h]$ is a deck transformation. Its conjugation on $\pi$ is the commutator of the identity automorphism with the automorphism induced by $\widetilde h$, hence is trivial. Centre-freeness makes this deck transformation the identity. This proves one inclusion in the displayed equality; descent proves the other. Finally $j(\id)=\id$, so a based loop has the same endpoint after lifting.
\end{proof}

This comparison does not supply a contraction inside the right-hand centralizer. Choosing a hyperbolic metric and interpolating from $j(f)(x)$ to $x$ along geodesics has two unproved properties: the arbitrary lift group $\widetilde H$ need not preserve these geodesics, and an interpolated map need not be a homeomorphism. The first prevents asserting equivariance; the second prevents asserting that the interpolation lies in $j(A^0)$. Neither property follows from the displayed centralizer equality. The universal vanishing of $\pi_1(C_{A^0}(H))$, and the requested counterexample to it, therefore remain unresolved by this argument.

\paragraph{Current orbitwise reduction and the remaining compatibility problem.}
C.26 settles the globally transitive centralizer case on a closed
hyperbolic surface and also gives an explicit closed subgroup with an
ambient-null but intrinsically nontrivial loop; its latter example is
not asserted to be a centralizer. For the canonical plane lift above,
C.28 identifies the orbits of the group generated by the loop:
points, intrinsic lines, circles, and open surface orbits.
C.27 gives an order-natural contraction on each intrinsic line.
C.29 proves that every circle restriction has winding zero by the
actual plane-loop evaluation map, and contracts that restriction
in its circle centralizer.
C.31 supplies inverse and track control, and C.34 proves joint
continuity of the contractions and their inverses on the entire
closed union $Y$ of the nonopen orbits. This is a based contraction
through homeomorphisms of $Y$, commuting with every homeomorphism
that commutes with the original loop. In particular C.34 settles
the hyperbolic centralizer loop whenever its lifted generated
group has no open orbit.

The open orbits have also been treated intrinsically. C.30 gives
the exact restricted operator group, including its restriction
kernel, and proves that its closed image acts freely and properly.
C.32 handles a discrete image, and C.33 treats every locally
compact closed image supplied by C.30. In each case the restricted
loop has a based null-homotopy inside the centralizer of the full
restricted operator group. These results give contractions on
all individual open orbits and on the whole nonopen locus.
The remaining step is to choose these contractions compatibly
across the frontiers of the open components, with joint
continuity of the resulting homeomorphisms and their inverses.

C.35 proves that this step cannot require extending the existing
median prescription unchanged. Its explicit plane action has
the boundary restriction map
\[
 \R^2\longrightarrow\R^3,\qquad
                  (A,B)\longmapsto(A,B,A-B).
\]
At the displayed parameter the prescribed three edge medians
are $(1,0,0)$, which lies outside its image because
$r_1-r_2-r_3=1$. Thus that prescribed homeomorphism of $Y$
has no equivariant extension over an adjacent open orbit.
The original loop in the same example has the explicit global
equivariant contraction given in (C.35.5).
C.36 constructs compatible choices on closures of open orbits
whose closed proper operator image has dimension two, using the
actual open quotient onto their simply transitive Lie
centralizers. Scaling the logarithm of each based universal-cover
lift respects the exact common boundary action. C.37 supplies a
uniform three-word estimate and proves joint continuity, including
inverses, on
\[
 X_2=\overline{\bigcup_{\dim Q_O=2}O},
\]
including every infinite accumulation of those open orbits.
Its bounded residual equations determine the unique logarithmic
limit on line orbits. The actual null residual word isotopies
and the two-point configuration map determine the real lifted
limit on circle orbits, including its winding.

C.49--C.50 subsequently enlarge this closed set to $X_*$,
including all open orbits with at least three actual boundary
quotient foliations. C.51 proves whole-closure extension for
the remaining open orbits with compact operator quotient,
using the median frontier prescription and finite pasting.
C.52 makes these extensions exactly relative to $X_*$.
C.53 enlarges the exponential region to $X^\dagger$ and
proves continuity across all accumulations of the operator
translates of finitely many remaining cocompact representatives.
C.54--C.56 enlarge the treated region through all
three-dimensional diagonal affine source groups, with
uniform accumulation control, and replace the compact
operator quotient condition by compactness of the full
local centralizer quotient. C.57 gives the exact plane,
open-cylinder, line and open-M\"obius quotient alternatives
in the remaining noncocompact two-family cases.
The outstanding problem includes untreated noncocompact local
quotients, infinitely many $H$-inequivalent remaining
representatives, and extension through exterior regions
consisting of nonopen orbits. The hybrid target prescription
has been proved continuous on the treated closed sets;
its transverse continuity on all of the nonopen locus has
not been proved. C.35 requires allowing changes to the
original median prescription. These exact remaining
requirements, and the general kernel in (C.12.1), remain
unresolved here.


\subsection{Configuration invariants and drilling}

\begin{claim}{C.13}{Proved}{The invariant braid group and nonabelian cohomology}
Let a finite group act faithfully and orientation preservingly on a
hyperbolic $U$, and let $R=U^!$ be its finite set of points with nontrivial
stabilizer. Put $V=U\setminus R$,
$B=\{f\in A^0:f|_R=\id\}$, and $B^0=\operatorname{Homeo}(V)^0$.
With the action by conjugation, we have
\[
 (B/B^0)^\Gamma=1,\qquad H^1(\Gamma,B^0)=1,
 \qquad (A^0)^\Gamma=B^\Gamma\cong
 \operatorname{Homeo}(V/\Gamma)^0.
\]
Here the quotient group identifies canonically with
$\pi_1(\operatorname{Conf}_R(U),R)$ by evaluation, and its invariant
subgroup is the $H^0$ asked for on scan 285.
\end{claim}
\begin{proof}
Evaluation $A^0\to\operatorname{Conf}_R(U)$ is a locally trivial
principal $B$-bundle, by the local extension theorem for embeddings of
a finite set. Its total space is contractible. Its homotopy exact
sequence gives $\pi_0B\cong\pi_1\operatorname{Conf}_R(U)$, with the
specified base configuration, and its identity component is the
identity component after retaining the new punctures. The construction
is equivariant for $\Gamma$, so this is the stated action on the braid
group.

An invariant component $b\in B/B^0$ is a mapping class of $V$ in the
kernel of filling $R$ that centralizes the action of $\Gamma$ on $V$.
By C.6 and its transporter version, choose a representative $f$ on $V$
commuting with $\Gamma$. It extends across the punctures and induces the
identity mapping class on $U$. Thus $f\in(A^0)^\Gamma$, and C.7 joins
it to the identity by an equivariant isotopy on $U$. Such an isotopy fixes
each point of $R$: it preserves each finite critical-point set and
starts with the identity permutation. Restriction to $V$ proves $b=1$.

For a nonabelian $1$-cocycle $a_\gamma\in B^0$, the maps
$v(\gamma)=a_\gamma u(\gamma)$ form another finite action on $V$.
The cocycle equation is exactly the homomorphism equation for $v$.
Both actions induce the same homomorphism into the mapping class group
of $V$. By C.6 they are conjugate by an element $b\in B^0$; equivalently
$a_\gamma=b u(\gamma)b^{-1}u(\gamma)^{-1}$. Hence the cocycle is a
coboundary and $H^1(\Gamma,B^0)=1$.

Finally C.7 gives connectedness. Descent to the quotient with its branch
points retained is precisely descent to $V/\Gamma$ with its punctures
retained; the identity lift construction in C.7 is the displayed
isomorphism. In particular the unprimed $V$ in the transcription of the
identity-component line on scan 286 must be compared with the quotient
$V'$ of scans 282--285. The two homeomorphism groups need not be
isomorphic; the descent map identifies the quotient one.
\end{proof}

\begin{claim}{C.14}{Refuted}{Finite critical configurations do not include mirrors}
The finite-set configuration formulation of C.13 does not extend by
replacing an orientation-preserving action by an arbitrary action and
keeping $R=U^!$ equal to all points with nontrivial stabilizer. A reflection
has a one-dimensional fixed locus, and deleting finitely many points
cannot make that action free.
\end{claim}
\begin{proof}
Take complex conjugation on
$U=\mathbf P^1(\C)\setminus\{0,1,\infty\}$. It acts faithfully on this
hyperbolic surface. Its fixed set is
$\mathbf P^1(\R)\setminus\{0,1,\infty\}$, a union of three open arcs.
Every point of these arcs has nontrivial stabilizer. A finite deletion
leaves points in each arc. At any remaining fixed point, the lift of
conjugation fixing a selected universal-cover point is an element of
order two in the lift extension. Therefore the drilled extension still
has torsion. This explicit counterexample refutes both the finiteness
premise and a torsion-free finite-drilling assertion for such an action.
The corrected geometry retains mirror boundary in the quotient or cuts
along the mirror locus; it is not the operation of adding finitely many
puncture loop classes. The equations C.13 were stated on scans 282--286
after the explicit restriction $\Gamma=\Gamma^+$.
\end{proof}

\begin{claim}{C.15}{Proved}{The oriented drilling equivalence and the bijection}
Fix a finite group $\Gamma$ and faithful orientation-preserving outer
actions on discrete hyperbolic groups with loops. Let $\Phi(\pi,\Gamma)$
be the set of $\pi$-conjugacy classes of nontrivial maximal finite
subgroups of the canonical extension. Filling a specified
$\Gamma$-stable set $J$ of punctures with nontrivial stabilizers induces
an equivalence between the category of such drilled actions and the
category of filled actions with a specified subset
$J\subset\Phi(\pi,\Gamma)$. In particular
\[
 \Phi(\pi',\Gamma)\amalg J\xrightarrow{\ \sim\ }\Phi(\pi,\Gamma)
\]
is a canonical $\Gamma$-equivariant bijection. Drilling all of $\Phi$
produces a torsion-free extension, uniquely up to the unique exterior
isomorphism over the specified filled object.
\end{claim}
\begin{proof}
Realize the action by C.5. A finite nontrivial subgroup of the lift group
acts on $\mathbb H^2$ by orientation-preserving isometries and has a
unique fixed point. Its stabilizer at that point is the unique maximal
finite subgroup containing it. Conversely every point of $U$ with
nontrivial stabilizer gives that subgroup after choosing a lift.
Changing the lift conjugates by $\pi$. This identifies $\Phi$ with the
finite critical-point set of the realized action, including its
$\Gamma$ action. It also proves the assertion about maximality, rather
than assuming it for arbitrary finite subgroups containing reflections.

Deleting the points of $J$ gives the drilled surface. Remaining critical
points and deleted points partition the original critical-point set,
which proves the displayed bijection. The group map is the epimorphism
$\pi'\to\pi$ obtained by killing the normal closure of exactly the new
peripheral subgroups. Local punctured discs show that its induced map
on lift groups sends a cusp stabilizer to the corresponding finite
point stabilizer. Thus the geometric bijection is the algebraic map
of scan 290, with the same conjugacy ambiguities.

To prove fullness, realize an equivariant outer isomorphism of the
filled peripheral groups by an equivariant homeomorphism, using C.6's
transporter conclusion. It preserves the selected subset of $\Phi$ by
the preceding construction. Its restriction to the complements gives
the desired drilled outer isomorphism. To prove faithfulness, represent
a drilled equivariant outer automorphism in the kernel by an equivariant
homeomorphism of the drilled surface and extend it across the new
punctures. It is isotopic to the identity on the filled surface. By C.7
there is an equivariant identity isotopy there. Every point of $J$ stays
fixed throughout, since it is a member of the finite critical set, so
the isotopy restricts to the drilled surface. The kernel is therefore
trivial. These two arguments also prove uniqueness over a fixed filled
isomorphism, not merely uniqueness of isomorphism classes.

After all critical points are removed, the finite action is free. Its
lift extension acts freely on $\mathbb H^2$ and is torsion-free, since
a nontrivial finite-order isometry fixes a point. The operations of
deleting and filling are inverse by the just-proved full faithfulness.
The only use of geometric realization in this proof is C.5 together
with the equivariant homeomorphism and isotopy results C.6--C.7; the
comparison map on groups and the kernel argument have been specified.
\end{proof}

\begin{claim}{C.16}{Refuted}{Nonfaithful outer actions obstruct torsion-free drilling}
The unrestricted existence claim on scans 286--288 fails for the
trivial outer action of $C_2$ on the group of type $(0,3)$. The faithful,
orientation-preserving version is C.15.
\end{claim}
\begin{proof}
Suppose a finite-puncture drilling of that action had torsion-free
extension. Realize its finite outer action on the drilled hyperbolic
surface by C.5. Filling the specified punctures extends these
homeomorphisms to a finite action on the filled three-punctured sphere.
The induced mapping class of every element is trivial. The finite-order
rigidity argument of C.5 makes every filled homeomorphism the identity.
Its restriction to the dense drilled surface is therefore the identity.
The upstairs outer action is trivial as well. For a centre-free kernel
its canonical extension is $\pi'\times C_2$, which contains
$\{1\}\times C_2$. This contradicts torsion-freeness. The argument proves
the same obstruction for every nonfaithful outer action after restricting
to a nontrivial element of its finite kernel.
\end{proof}

\begin{claim}{C.17}{Proved}{Torsion-free finite extensions with peripheral data}
Let $\pi$ be a discrete hyperbolic surface group, possibly punctured,
and let $1\to\pi\to E\to\Gamma\to1$ be a torsion-free extension by a
finite group whose outer action preserves the loop structure. Then $E$
is the fundamental group of a finite-type surface, possibly
nonorientable. Its peripheral subgroups are the centralizers in $E$ of
the peripheral subgroups of $\pi$. With their conjugacy classes it is a
surface group with loops in this extended sense. The associated compact
surface and boundary groups form a Poincar\'e duality pair of dimension
two with its orientation local system.
\end{claim}
\begin{proof}
The centre-free extension is the pullback of
$\Aut(\pi)\to\Out(\pi)$. The centralizer of $\pi$ in $E$ maps
isomorphically to the kernel of the outer action: a representative
inducing an inner automorphism can be multiplied by the inverse inner
element to centralize $\pi$, and this representative is unique because
$Z(\pi)=1$. This centralizer is finite, so torsion-freeness forces it
to be trivial. Hence the outer action is faithful.

Apply C.5 and identify $E$ with the group of lifts of the realized
action, by uniqueness of the centre-free extension. A nontrivial
stabilizer of a point gives a finite subgroup of $E$, so the action on
$U$ is free. Its quotient $W=U/\Gamma$ is an ordinary surface and
$\mathbb H^2\to W$ is its universal covering. This gives the required
isomorphism $E\cong\pi_1(W)$, respecting the specified subgroup
$\pi$ and quotient $\Gamma$.

A stabilizer of an end of $U$ cannot contain an orientation-reversing
operator. The local finite-group model at its compactifying point would
then contain a reflection and fixed points arbitrarily close to that
point, contrary to freeness on $U$. Its local action is therefore cyclic
rotation. Consequently an upstairs peripheral subgroup is a finite-index
subgroup of the downstairs peripheral infinite cyclic group. In the
hyperbolic universal cover its centralizer is exactly that full cusp
group. These are all the cusps, since every end downstairs has a
preimage upstairs. Thus the stated peripheral classes, including their
full subgroups rather than merely inclusions, are obtained.

Choose disjoint cusp neighbourhoods in $W$ and truncate them to obtain
a compact surface $M$ with boundary. The inclusions of boundary circles
induce precisely these cyclic groups. Lift a triangulation to the
contractible universal cover. The cellular cap-product map with the
relative orientation class gives the chain equivalence between
cochains and complementary relative chains; the local sign is the
orientation character of $W$. This is the Poincar\'e--Lefschetz duality
pair, and its group-pair chain model is the relative cellular complex
because both the universal cover and the lifted boundary components
are contractible. This proves the duality-pair assertion rather than
inferring it from virtual freeness alone. The exact geometric input is
the finite realization theorem; a separate proof from abstract duality
axioms is not substituted for this argument.
\end{proof}

\begin{claim}{C.18}{Proved}{The discrete oriented realization criterion}
For a faithful orientation-preserving finite outer action on a discrete
hyperbolic group with loops, the criterion of scan 294 is valid with
the word admissible meaning the full-faithfulness and uniqueness
conditions of scan 293. The drilling in question is at the canonical
set $\Phi(\pi,\Gamma)$, and its extension carries the quotient loop
structure described in C.17.
\end{claim}
\begin{proof}
A realized action can be drilled at its critical points. C.15 proves
that the resulting action is admissible and its extension torsion-free;
C.17 identifies that extension with a group with loops. This verifies
every requirement of the stated criterion.

Conversely let $\pi'\to\pi$ be the stipulated filling map and let
$E'$ be its extension, with the group-with-loops structure required in
the criterion. Realize $E'$ by a compact surface with its boundary
circles, then take its covering corresponding to $\pi'$. The comparison
of peripheral subgroups in the definition ensures that this cover has
exactly the prescribed upstairs loop classes. Its deck group is
$\Gamma$, acting freely before filling. The cyclic end stabilizers
extend across each new puncture by the local map $z\mapsto z^m$.
Filling the selected ends thus gives a $\Gamma$-action on the surface
whose peripheral group is $\pi$. The canonical map of lift extensions
identifies its outer action with the original one. Admissibility makes
this construction independent, up to the uniquely specified exterior
isomorphism, of the chosen drilled object. This proves the converse
by the typed covering and filling maps. The mirror counterexample
C.14 explains the necessary restriction of this formulation.
\end{proof}

\subsection{Special homeomorphisms: the topology and the extension class}

\begin{claim}{C.19}{Refuted}{Escaping support invalidates the compact-open version}
Let $SA\subset A$ consist of homeomorphisms equal to the identity on
some neighbourhood of every puncture, and give it the subspace
compact-open topology. The canonical special-group comparison of
scan 304 does not identify $\pi_0(SA)$ with the boundary-framed mapping
class group. A peripheral boundary twist is connected to the identity
inside this $SA$. In particular, the annulus conclusion
$\pi_0(SA)=\Z$ on scan 304 is false in this topology.
\end{claim}
\begin{proof}
Use a coordinate $z=re^{i\theta}$ on a disc about a puncture at $0$.
Choose a continuous function $\eta:[0,\infty)\to[0,1]$ equal to $0$ on
$[0,1/3]$ and to $1$ on $[2/3,\infty)$. For $0<s\leq1$ define
\[
 h_s(re^{i\theta})=r\exp\bigl(i\theta+2\pi i\eta(r/s)\bigr),
 \qquad h_0=\id,
\]
and extend by the identity outside the coordinate disc. Each $h_s$ is
identity near the puncture and outside the annulus
$s/3<r<2s/3$. It is a boundary twist when viewed on a fixed compact
surface obtained by truncating the puncture. Nevertheless
$\sup_z|h_s(z)-z|\leq4s/3$, and the same bound holds for its inverse.
Thus $s\mapsto h_s$ is a continuous path in the stated subspace topology
from the identity to that twist.

For the annulus every compactly supported orientation-preserving
homeomorphism fixes both ends and is represented on a compact annulus
by an integer twist; cutting along a spanning arc and applying the
Alexander isotopy of a disc proves this integer classification. The
isotopy to the corresponding integer twist can be chosen in that fixed
compact annulus. The path above then kills that twist in $SA$. Hence
$SA$ is path connected in the compact-open topology, proving the
explicit contradiction. The programme checks the formula and its
uniform bound; the bound itself proves convergence for the full space,
not just on a finite sample.

The fixed-boundary formulation below retains a single compact surface
throughout the isotopy. That is the required correction. No change of
topology is silently made in this refutation.
\end{proof}

\begin{claim}{C.20}{Proved}{The fixed-boundary repair and its class}
Let $M$ be an oriented compact surface obtained by replacing the
$\nu>0$ punctures of a hyperbolic type by boundary circles. Put
$A_\partial=\operatorname{Homeo}^+(M\ {\rm rel}\ \partial M)$.
Then $A_\partial^0$ is contractible,
$\pi_0(A_\partial)\cong S\Gamma^{!+}_{g,\nu}$ for the framed group
of E.1--E.2, and there is a central extension
\[
 1\longrightarrow\Z^\nu\longrightarrow\pi_0(A_\partial)
 \longrightarrow\Gamma^{!+}_{g,\nu}\longrightarrow1.
\]
The $i$th kernel generator is the positive boundary twist. Its
extension class is the vector of Euler classes of the oriented tangent
circle lines at the marked points, with the sign fixed by that generator.
For the annulus the kernel is the quotient of $\Z^2$ by the diagonal
rotation relation; for the once-punctured sphere it is zero.
\end{claim}
\begin{proof}
Allow boundary rotations first, retaining each boundary component
setwise. Restriction to the boundary is a locally trivial fibration
onto $\prod_i\operatorname{Homeo}^+(S^1)$: extend a small boundary
isotopy across collars. Its fibre is $A_\partial$.
Each boundary homeomorphism group deformation retracts onto rotations
by lifting to increasing maps of $\R$ and linearly interpolating the
periodic displacement. Thus the base has homotopy type $(S^1)^\nu$.
Capping each boundary by a disc with its centre marked identifies the
free-boundary mapping classes with the pure oriented punctured mapping
classes. The same collar-extension construction gives the homotopy
comparison on identity components.

The identity component of the free-boundary group is contractible in
a hyperbolic type by Yagasaki's compact-pair theorem. The homotopy
sequence therefore injects $\pi_1((S^1)^\nu)=\Z^\nu$ into the
components of the fibre and identifies the quotient as stated. Lifting
one rotation in the $i$th collar gives exactly the boundary twist.
Its conjugation action is trivial because the boundary circles are
retained individually and their orientations are fixed. Choose a point
on each boundary circle. A boundary-fixing homeomorphism acts on the
framed peripheral diagram of E.2, and an isotopy relative to the
boundary gives its prescribed natural isomorphism. This defines
$\pi_0(A_\partial)\to S\Gamma^{!+}_{g,\nu}$. On the quotient it is
the peripheral Dehn--Nielsen--Baer isomorphism. A positive boundary
twist changes the path from the generic basepoint to that boundary
point by one peripheral turn; hence the induced map on kernels sends
the $i$th boundary generator to the $i$th generator of $T^I$.
It is an isomorphism on both kernel and quotient, so the two short
exact sequences give the asserted isomorphism of extensions.
Contractibility of $A_\partial^0$ follows directly from Lemma 2.1(ii),
with the nonempty fixed boundary. The connecting homomorphism is consequently the
fundamental-group extension of the bundle of tangent-circle framings
at the marked points. Obstruction to a section of that circle bundle
is its Euler class, proving the class identification with its sign
convention. In complex notation this is the tangent-line first Chern
class; the usual cotangent $\psi_i$ has its opposite sign.

For the annulus the identity component of the free-boundary group is
$S^1$. Its fundamental-group map to the two boundary-rotation groups
is the diagonal when both circles use the same angular coordinate,
or the antidiagonal when both use boundary orientations. The cokernel
is $\Z$. For the disc this map is an isomorphism $\Z\to\Z$, so the
cokernel is zero. These are the exact boundary models underlying the
two abelian computations on scans 303--304.
\end{proof}

\begin{claim}{C.21}{Proved}{The four small extension classes}
With the positive boundary-twist convention of C.20, the special
fixed-boundary extension classes for $(0,3)$ and $(0,4)$ are zero. For
$(1,1)$ the class is a generator $e$ of
$H^2(\operatorname{SL}_2(\Z);\Z)=\Z/12$. For $(1,2)$ the vector class
is $(q^*e,q^*e)$, where
$q:\Gamma^{!+}_{1,2}\to\operatorname{SL}_2(\Z)$ forgets the second
marked point; its order is $12$.
\end{claim}
\begin{proof}
The pure oriented mapping class group of a three-punctured sphere is
trivial. Filling one puncture of a four-punctured sphere gives the
point-pushing group $F_2$, since the target pure group is trivial and
the three-punctured surface is hyperbolic. A graph is a classifying space
for $F_2$, so $H^2(F_2;\Z^4)=0$. This proves the first two class
computations, including the noncanonical splitting in type $(0,4)$.

The integral matrices
\[
 S=\begin{pmatrix}0&-1\\1&0\end{pmatrix},\qquad
 R=\begin{pmatrix}0&-1\\1&1\end{pmatrix}
\]
have orders $4$ and $6$ and satisfy $S^2=R^3=-I$.
Their action on the Farey tessellation gives
$\operatorname{PSL}_2(\Z)=C_2*C_3$: the dual bipartite tree has one
edge as fundamental domain, vertex stabilizers $C_2,C_3$, and trivial
edge stabilizer. Connectivity and the absence of circuits follow from
the Euclidean algorithm for adjacent reduced fractions. Lifting the
central element $-I$ gives
$\operatorname{SL}_2(\Z)=C_4*_{C_2}C_6$. The cohomology sequence of
this tree of groups identifies $H^2$ with
\[
 \{(r,s)\in\Z/4\oplus\Z/6:r\equiv s\pmod2\}:
\]
the preceding $H^1(C_2;\Z)$ and $H^1(C_4;\Z)\oplus
H^1(C_6;\Z)$ vanish, while the restrictions in degree two reduce
modulo $2$. This group has twelve elements and is cyclic.
The tangent-circle bundle restricts on the square torus fixed by $S$
to the rotation representation by $\pi/2$, and on the hexagonal torus
fixed by $R$ to the representation by $\pi/3$. Lifting these circle
rotations to $\R$ gives fourth and sixth powers equal to one positive
full turn. By the convention of C.20 the Euler class therefore
restricts to $(1,1)$, an element of order $12$. This proves its generator
assertion without an unproved presentation for the boundary mapping
class group. The supplied code verifies the matrix orders and the
finite abelian-group calculation.

In a complex torus translation canonically identifies the tangent
spaces at any two points. These identifications vary with the marked
complex torus and commute with all isomorphisms. Hence the two tangent
circle bundles on the two-pointed moduli orbifold are isomorphic, and
both are pulled back from the one-pointed tangent bundle. This proves
the vector formula. It has order dividing $12$. A square torus with
marked points $0$ and $(1+i)/2$ admits the order-four rotation preserving
both; a hexagonal torus admits an order-three rotation with at least
two marked fixed points. The restrictions of the class to these cyclic
subgroups are generators, as can be seen from their nontrivial local
rotations, or from C.22 below. Its order is divisible by $4$ and $3$,
and thus is exactly $12$. The existence of the displayed fixed points
is also verified in the integral torus matrices supplied with the code.
\end{proof}

\begin{claim}{C.22}{Proved}{Restriction over a finite puncture stabilizer}
For a hyperbolic type, a nontrivial finite subgroup
$G\subset\Gamma_i^+$ is cyclic. The extension by the $i$th copy of
$T=\Z$ obtained from the special group is nonsplit over $G$; its
preimage is infinite cyclic and the kernel is the subgroup of index
$|G|$.
\end{claim}
\begin{proof}
Realize $G$ by C.5. It fixes the marked point $i$ and preserves
orientation. The finite local action is a faithful rotation action:
an element trivial on a small disc is the identity as a periodic
surface homeomorphism, by the local classification and connectedness.
Thus $G$ embeds in the circle rotations and is cyclic. Delete that
point and lift a small punctured quotient disc to its universal cover.
Its group is $\Z$, while the upstairs punctured disc has subgroup
$|G|\Z$. The exact normalizer interpretation proved in E.3 identifies
this with the extension pushed out by the $i$th projection. An infinite
cyclic group contains no nontrivial finite subgroup, so the extension
cannot split. This proves the remark on scan 306 with its stated
orientation and puncture-stabilizer hypotheses.
\end{proof}


\begin{claim}{C.23}{Proved}{An infinite-group failure of existence}
For every closed orientable surface $X_g$ of genus $g\geq2$, the
identity homomorphism from the discrete group
$\Gamma=\operatorname{Mod}^+(X_g)$ to its subgroup in $\pi_0A_g$
has no homomorphic lift to $A_g$. Thus level $3$ fails for an infinite
operator group already on each such closed hyperbolic surface.
\end{claim}
\begin{proof}
Markovi\'c--\v Sari\'c's Theorem 1.1
\cite{C_MarkovicSaric2008} states that the projection
$\operatorname{Homeo}^+(X_g)\to\operatorname{Mod}^+(X_g)$ has no group
section for $g\geq2$. The theorem is about homeomorphisms, not merely
diffeomorphisms; its first-page statement and definitions have been
checked in the original paper. A lift of the displayed homomorphism
would take every element to an orientation-preserving homeomorphism,
since orientation is invariant under isotopy. Its composite with that
projection would be the identity, giving the forbidden section. This
is the exact reduction to the source theorem. The group is infinite:
a twist about a nonseparating simple curve acts on $H_1(X_g;\Z)$ by
an integral transvection, whose distinct powers are distinct.
\end{proof}

\subsection{The integral sphere comparison}

\begin{claim}{C.24}{Refuted}{The sphere comparison already fails for an order-three rotation}
Let $A=\operatorname{Homeo}(S^2)$ with its compact-open topology, and let
$\Gamma=C_3$ act on $S^2$ by rotation through $2\pi/3$ about a fixed axis.
Give $A$ the resulting conjugation action and base both fixed-point spaces
at the identity. The actual comparison has
\[
 \pi_1(A^\Gamma,1)=\Z,\qquad
 \pi_1(A^{h\Gamma},1)=\Z_3\times\Z/2,
 \qquad n\longmapsto(n,n\bmod2).
\]
Consequently the requested integral homotopy equivalence fails in type
$(0,0)$. This failure concerns the actual homeomorphism mapping space;
it is independent of the earlier loss of the sphere's $\pi_2$ in its
fundamental groupoid. For every odd prime $p$, the same construction with
$C_p$ gives $\Z_p\times\Z/2$ and the same comparison formula with $p$
in place of $3$.
\end{claim}
\begin{proof}
We give the comparison maps and compute their effect, including the integral
rather than only completed homotopy group. Write $G=SO(3)$ and
$\rho:C_p\hookrightarrow G$ for the rotation subgroup, with $p$ odd.
The inclusion $G\hookrightarrow A^+=\operatorname{Homeo}^+(S^2)$ is
$C_p$-equivariant and a weak equivalence. Indeed evaluation on three
ordered distinct points is a locally trivial fibration onto
$\operatorname{Conf}_3(S^2)$, by isotopy extension. Its fibre is the
pure homeomorphism group of the three-punctured sphere. Its component
group is trivial, and its identity component is contractible by the
surface homeomorphism-group theorem used in C.7. The unique M\"obius
transformation carrying one ordered triple to another identifies this
configuration space with $PSL_2(\C)$; polar decomposition retracts this
group onto $SO(3)$. The resulting evaluation diagram proves the stated
weak equivalence for the actual inclusion. Homotopy fixed points preserve
this objectwise weak equivalence, after applying singular complexes. A
path through the identity in $A^{hC_p}$ stays in $A^+$, because $EC_p$
is connected and orientation takes values in the discrete group $C_2$.
It therefore suffices to calculate $(G^{hC_p},1)$.

The identity component of the actual centralizer fixes the two poles
individually. Quotient by $z\mapsto z^p$ identifies its descent map to
$\operatorname{Homeo}^+(S^2;0,\infty)^0$ as a covering of degree $p$.
Every map of the quotient fixing both poles lifts: on the complementary
annulus its action on the meridian is $+1$, so it preserves the subgroup
$p\Z$ of the covering group. Each lift commutes with the deck group,
since the induced automorphism of $\Z/p$ is $+1$; it extends uniquely
over the two poles. The usual local lifting neighbourhoods establish
that descent is a covering in the compact-open topology. Its total space
is connected: the quotient group is connected, and all deck lifts lie
on the connected circle of axis rotations. The annulus calculation of
C.10 identifies the base with a space of homotopy type $S^1$. The
rotation circle maps to its quotient rotation circle by $z\mapsto z^p$;
its inclusion in the total centralizer therefore induces an isomorphism
on every homotopy group. In particular the actual fundamental group is
$\Z$, with generator one complete axis rotation.

Set $Y=BG=BSO(3)$. The adjoint bundle
$EC_p\times_{C_p}G\to BC_p$, with action through $\rho$, has section
space $G^{hC_p}$. It is also the bundle of gauge automorphisms of the
principal $G$-bundle classified by $B\rho$. The universal bundle gives
the natural weak equivalence
\[
 G^{hC_p}\simeq
 \Omega_{B\rho}\operatorname{Map}(BC_p,Y).
\]
One way to verify the map is to send a gauge automorphism to the loop
of its classifying map obtained by gluing the bundle over the two ends
of $BC_p\times[0,1]$. Conversely a loop of classifying maps gives a
bundle over this cylinder with its specified endpoint identifications;
trivialization in the interval direction gives the gauge automorphism.
The spaces of such trivializations are contractible. This construction
is natural for the inclusion $C_G(\rho(C_p))\hookrightarrow G^{hC_p}$.

We now compute $\pi_2$ of this mapping space using the arithmetic square
of the simply connected finite-type space $Y$. Let
\[
 \widehat Y=\prod_qY_q^{\wedge},\qquad
 V=\Bigl(\prod_q\Z_q\Bigr)\otimes\Q.
\]
Here $q$ runs over primes and superscript $\wedge$ means homotopical
$q$-completion, throughout this proof. We have
$Y_{\Q}\simeq K(\Q,4)$ and
$(\widehat Y)_{\Q}\simeq K(V,4)$. This follows from
$SO(3)\simeq\mathbb {RP}^3$: its rational homotopy is concentrated in
degree three, and each of its higher homotopy groups is finite. The
square
\[
 \begin{matrix}
 Y&\longrightarrow&\widehat Y\\
 \downarrow&&\downarrow\\
 K(\Q,4)&\longrightarrow&K(V,4)
 \end{matrix}
\]
is a homotopy pullback. For completeness its homotopy-group verification
reduces on each free summand to
\[
 0\longrightarrow\Z\longrightarrow
 \Q\oplus\prod_q\Z_q\longrightarrow V\longrightarrow0,
 \quad (a,(b_q))\longmapsto(b_q-a)_q.
\]
Surjectivity follows by clearing a common denominator and choosing its
finitely many required residues by the Chinese remainder theorem; the
kernel consists of rational numbers integral at every prime and is
$\Z$. On finite summands completion gives the product of their primary
parts and rationalization is zero. The long exact sequence of the
homotopy pullback and Whitehead's theorem prove the square, also giving
the usual arithmetic-square construction in the exact form needed here.

Apply $\operatorname{Map}(BC_p,-)$ to this homotopy pullback and take the
component containing $B\rho$. For every rational vector space $W$,
positive-degree cohomology $H^k(C_p;W)$ vanishes by averaging. The two
rational mapping spaces consequently have vanishing $\pi_2$ and
$\pi_3$. The homotopy-pullback exact sequence therefore gives
\[
 \pi_2\operatorname{Map}(BC_p,Y)_{B\rho}
 \cong\prod_q
 \pi_2\operatorname{Map}(BC_p,Y_q^{\wedge})_{B\rho}.
\]
For $q\ne p$, evaluation is a weak equivalence from this mapping-space
component to $Y_q^{\wedge}$. To check it, use the Postnikov tower:
positive-degree $C_p$-cohomology of every $q$-complete homotopy group
vanishes since multiplication by $p$ is invertible; at each tower stage
only constant maps contribute, and homotopy inverse limits preserve
the comparison. Thus only $q=2$ contributes to $\pi_2$ among these
primes, giving $\pi_2(Y_2^{\wedge})=\Z/2$.

At $q=p$ we use the exact finite-$p$-group mapping-space theorem:
for a finite $p$-group $P$, a compact Lie group $K$, and
$\alpha:P\to K$, the map induced by the centralizer identifies
\[
 \pi_n\operatorname{Map}(BP,(BK)_p^{\wedge})_{B\alpha}
  =\bigl(\pi_n B C_K(\alpha(P))\bigr)_p^{\wedge}
  \quad(n\ge2).
\]
This is the Dwyer--Zabrodsky--Notbohm formula, stated and derived in
\cite[Theorem 1.9 and its proof, using Theorem 3.2(iii) of their earlier
paper]{C_JMO1995}. Its hypotheses apply to the displayed finite cyclic
group and $SO(3)$. Since a nontrivial odd-order rotation has centralizer
exactly the axis circle $SO(2)$, the contribution is $\Z_p$.
This proves $\pi_1(G^{hC_p},1)=\Z_p\times\Z/2$.

Naturality of the centralizer map makes the $p$-component of the map
from the rotation circle the completion $\Z\to\Z_p$. At the prime
$2$ it is the map $\pi_1SO(2)\to\pi_1SO(3)$, namely reduction modulo
$2$: lifting a full rotation to unit quaternions gives the endpoint
$-1$. This proves the stated formula for the exact comparison.
For an explicit missing element at $p=3$, put
$\beta=\sum_{j\ge1}3^{j!}\in\Z_3$. Its base-three digits have
infinitely many zeros and infinitely many ones, whereas an ordinary
integer has eventually only zeros or eventually only twos. Hence
$(\beta,0)$ is outside the image. This also proves nonsurjectivity
without a cardinality argument; for general odd $p$ the same factorial
series has digits zero and one, rather than eventually zero or $p-1$. The explicit action, its finite order, the axis
centralizer and the finite reductions of this comparison are checked in
\texttt{code/verify\_sphere\_comparison.py}. The inverse-limit argument
above, rather than a finite computation, proves the nonsurjectivity.
\end{proof}

\begin{claim}{C.25}{Refuted}{Zero-level fidelity for a faithful infinite cyclic sphere action}
Let $\Gamma=\Z$ act faithfully on $S^2$ by an irrational axis rotation,
$u(1)(\theta,t)=(\theta+2\pi\alpha,t)$ with
$\alpha\in\R\setminus\Q$ and the usual collapse of the latitude
circles at $t=\pm1$. Let $M=A^\Gamma$ be its actual automorphism
space, based at the identity, and let $N$ be its weak automorphism
space in the categorical group $\Pi_1A$, where
$A=\operatorname{Homeo}(S^2)$. The induced transporter map has
\[
 \pi_1(M,1)=\Z\longrightarrow\pi_1(N,1)=\Z/2,
 \qquad n\longmapsto n\bmod2.
\]
Its kernel is $2\Z$, so level $0$ fails for this faithful action of an
infinite group. This is the exact fundamental-groupoid comparison;
C.24 separately retains the higher homotopy of $A$ when computing the
untruncated homotopy-fixed space. The hyperbolic zero-level question
recorded in C.12 has its stated, stronger surface hypothesis.
\end{claim}
\begin{proof}
Commutation with an irrational rotation implies commutation with every
axis rotation: its powers are dense in $S^1$, and composition and
inversion in the compact-open homeomorphism group are continuous.
An equivariant homeomorphism preserves the two poles as a set and
sends each other rotation orbit to another orbit. In the orientation-
preserving identity component it has exactly the form
\[
 f(\theta,t)=(\theta+a(t),b(t)),\quad
 a\in C((-1,1),\R/2\pi\Z),\quad
 b\in\operatorname{Homeo}^+([-1,1]).
\]
To prove that this parametrization is onto, equivariance forces the
angular function on each orbit to be translation by a unique $a(t)$;
the orbit quotient supplies $b$. Orientation forces $b$ to be increasing.
Conversely every such pair extends continuously at both poles, because
those circle orbits collapse there, and the same argument applies to
its inverse. Thus the parametrization is a homeomorphism of spaces
with the compact-open topologies. Continuity at a pole imposes no
limit on $a(t)$; on any compact subinterval of $(-1,1)$ it is exactly
ordinary compact-open convergence.

The homotopy $a_s(t)=a((1-s)t)$ contracts the angular-function factor
to its evaluation circle at $0$. The convex interpolation
$b_s(t)=(1-s)b(t)+st$ contracts the increasing-homeomorphism factor
to the identity: it remains strictly increasing and fixes both endpoints.
These homotopies are continuous in the stated topologies and extend
continuously to the sphere. The actual identity centralizer therefore
has homotopy type $S^1$, with its axis-rotation circle giving the
specified generator of $\pi_1M$.

In the weak transporter groupoid, an automorphism of its identity
object is a loop in $A$ fixed by the conjugation action of $\Gamma$.
This is exactly the groupoid calculation of C.9, now applied to
$\Pi_1A$ rather than asserting any higher truncation of $A$ itself.
By the evaluation and polar-decomposition comparison in C.24,
$\pi_1(A,1)=\pi_1SO(3)=\Z/2$. Conjugation by an axis rotation acts
trivially on this group: joining that rotation to the identity in
$SO(3)$ gives a based homotopy of its conjugation map to the identity.
Hence $\pi_1N=\Z/2$.

The transporter comparison sends a loop of equivariant homeomorphisms
to that same loop in $A$. Its effect is consequently the usual
$\pi_1SO(2)\to\pi_1SO(3)$, reduction modulo two. More explicitly,
a loop of two complete rotations has degree two under the evaluation
$a\mapsto a(0)$ in the actual centralizer, so it is nonzero there.
Its lift to the quaternion double cover $S^3\to SO(3)$ is the closed
loop $t\mapsto\cos(2\pi t)+\mathbf i\sin(2\pi t)$, which contracts
in the simply connected $S^3$. It is therefore killed by the weak
transporter map. This proves the exact nontrivial kernel and the
zero-level failure, with no inference from disconnectedness.
\end{proof}
